\documentclass[12pt]{jlreq}
\usepackage{amsmath, amssymb}

\newcommand{\C}{\mathbb{C}}
\newcommand{\R}{\mathbb{R}}
\newcommand{\Z}{\mathbb{Z}}
\newcommand{\Tr}{\operatorname{Tr}}
\newcommand{\Impart}{\operatorname{Im}}
\newcommand{\AF}{\operatorname{AF}}
\newcommand{\AX}{\operatorname{AX}}
\newcommand{\GL}{\operatorname{GL}}

\begin{document}

\(3\) 次元ユークリッド空間 \(\R^3\) の原点を \(O\) とする。

\begin{enumerate}
  \item 以下の \(3\) つの条件を満たす \(C^\infty\) 写像 \(F:\R^3\to\R^3\) は，ある定数 \(a,b,c\) を用いて
  \[
    F(x,y,z)=(ax,by,cz)
  \]
  と書けることを示せ。
  \begin{enumerate}
    \item[(i)] \(\R^3\) 内の \(x\) 軸に平行な任意の直線は，\(F\) により \(x\) 軸に平行な直線に写される。
    同様に，\(y\) 軸に平行な任意の直線は \(y\) 軸に平行な直線に，また \(z\) 軸に平行な任意の直線は \(z\) 軸に平行な直線に写される。
    \item[(ii)] \(F\) のヤコビ行列式は，\(0\) ではない定数値関数である。
    \item[(iii)] \(F(O)=O\).
  \end{enumerate}
  \item 以下の \(4\) つの条件を満たす \(C^\infty\) 写像 \(G:\R^3\to\R^3\) は線形写像に限られることを証明せよ。
  \begin{enumerate}
    \item[(o)] \(G\) は単射である。
    \item[(i)] \(G\) は \(\R^3\) 内の任意の平面を，\(\R^3\) 内の平面に写す。
    \item[(ii)] \(G\) のヤコビ行列式は，\(0\) ではない定数値関数である。
    \item[(iii)] \(G(O)=O\).
  \end{enumerate}
\end{enumerate}

\end{document}